% claim.tex -- AI4Math claim of record (AI-generated, 2026-09-29; instantiated
% from harness/templates/claim/claim.tex per SOP 08b).
% Machine-readable theorem anchors are the "% LEAN:" lines; scripts/paper-lint.sh
% and the claim audit check them. All Lean listings are verbatim, byte-identical
% contiguous excerpts of frozen artifacts of tasks/20260929-comb10-ladder-pend3:
%   statements : formalized/comb10-ladder-pend3.statement.txt (160 lines, LF,
%                sha256 47d54db7...; byte-identical to formalized/comb10-ladder-pend3.lean)
%   proof file : final/comb10-ladder-pend3-proved.lean (440 lines, CRLF on the
%                working disk; excerpts LF-normalised per the chorded-cycle
%                precedent, byte-exact original shipped in proofs/)
% Line ranges used: statement 37-89, 92-116, 119-159 (base defs and theorem
% heads); final 1-440 (whole file). The frozen statement file's literal "sorry"
% occurrences (the three theorem proof placeholders at lines 90/117/160, plus
% one mention in the head comment block at line 33) all lie outside the
% excerpt ranges, as required by the paper-lint listing gate. Verification:
% audit/listing-verify.txt (audit/inspect-tex.py --verify).
\documentclass[11pt]{article}

\usepackage[T1]{fontenc}
\usepackage[utf8]{inputenc}
\usepackage{amsmath,amssymb,amsthm}
\usepackage{graphicx}
\usepackage{hyperref}
\usepackage{xurl}
\setlength{\emergencystretch}{3em}
% lstlean.tex — Lean style for the listings package.
% Lineage: Jeremy Avigad (2015), adapted from Assia Mahboubi's SSR style;
% inspired by Olivier Verdier's unixode.sty for the Unicode literate table.
% No CTAN package or upstream repo exists; papers carry this file in their
% source package (cap set arXiv:1907.01449, sphere eversion arXiv:2210.07746).
% AI4Math adaptation (AI-generated, 2026-09-26): sorry/admit/sorryAx elevated
% to a dedicated error tier, matching the pipeline's 0-sorry constitution.
\usepackage{listings}
\usepackage{xcolor}

\definecolor{keywordcolor}{rgb}{0.7, 0.1, 0.1}   % red keywords
\definecolor{tacticcolor}{rgb}{0.0, 0.1, 0.6}    % blue tactics
\definecolor{commentcolor}{rgb}{0.4, 0.4, 0.4}   % grey comments
\definecolor{symbolcolor}{rgb}{0.0, 0.1, 0.6}    % blue symbols
\definecolor{sortcolor}{rgb}{0.1, 0.5, 0.1}      % green sorts
\definecolor{errorcolor}{rgb}{0.6, 0, 0}         % dark red: sorry/admit
\definecolor{stringcolor}{rgb}{0.5, 0.1, 0.5}    % purple strings

\lstdefinelanguage{lean}{
  % Anything between $ becomes LaTeX math mode
  mathescape=false,
  % Comments may or not include Latex commands
  texcl=false,
  % keywords, list taken from lean-syntax.txt
  morekeywords=[1]{import, prelude, protected, private, noncomputable,
    definition, meta, renaming, hiding, parameter, parameters, begin,
    constant, constants, lemma, variable, variables, theory, print,
    theorem, example, open, section, namespace, universe, universes, end,
    modifier, modifiers, using, export, exposing, axiom, inductive,
    coinductive, with, structure, class, instance, attribute, local, where,
    by, have, show, let, in, if, then, else, match, do, for, fun, assume,
    from, take, calc, include, omit, infix, infixl, infixr, notation,
    macro, syntax, elab, set_option, initialize, deriving, opaque, abbrev,
    mutual, partial, scoped, termination_by, decreasing_by},
  % Sorts
  morekeywords=[2]{Sort, Type, Prop},
  % tactics, list taken from lean-syntax.txt (tactic keywords)
  morekeywords=[3]{intro, intros, apply, exact, refine, rfl, simp, simp_all,
    simpa, simp_rw, rw, rewrite, erewrite, ring, ring_nf, omega, decide,
    norm_num, norm_cast, induction, cases, rcases, obtain, constructor,
    ext, funext, conv, congr, field_simp, linarith, nlinarith, positivity,
    tauto, trivial, aesop, use, by_cases, by_contra, push_neg, contrapose,
    symm, trans, assumption, contradiction, exfalso, specialize,
    generalize, revert, clear, rename_i, subst, unfold, delta, change,
    convert, split_ifs, interval_cases, fin_cases, zify, gcongr,
    nth_rewrite, suffices, generalizing, at, only, native_decide},
  % warnings - constitution tier: must never survive into a final proof
  morekeywords=[4]{sorry, admit, sorryAx},
  literate=
    {α}{{\ensuremath{\alpha}}}1
    {β}{{\ensuremath{\beta}}}1
    {γ}{{\ensuremath{\gamma}}}1
    {δ}{{\ensuremath{\delta}}}1
    {ε}{{\ensuremath{\varepsilon}}}1
    {ζ}{{\ensuremath{\zeta}}}1
    {η}{{\ensuremath{\eta}}}1
    {θ}{{\ensuremath{\theta}}}1
    {ι}{{\ensuremath{\iota}}}1
    {κ}{{\ensuremath{\kappa}}}1
    {λ}{{\ensuremath{\lambda}}}1
    {μ}{{\ensuremath{\mu}}}1
    {ν}{{\ensuremath{\nu}}}1
    {ξ}{{\ensuremath{\xi}}}1
    {π}{{\ensuremath{\pi}}}1
    {ρ}{{\ensuremath{\rho}}}1
    {σ}{{\ensuremath{\sigma}}}1
    {τ}{{\ensuremath{\tau}}}1
    {υ}{{\ensuremath{\upsilon}}}1
    {φ}{{\ensuremath{\varphi}}}1
    {χ}{{\ensuremath{\chi}}}1
    {ψ}{{\ensuremath{\psi}}}1
    {ω}{{\ensuremath{\omega}}}1
    {Γ}{{\ensuremath{\Gamma}}}1
    {Δ}{{\ensuremath{\Delta}}}1
    {Θ}{{\ensuremath{\Theta}}}1
    {Λ}{{\ensuremath{\Lambda}}}1
    {Ξ}{{\ensuremath{\Xi}}}1
    {Π}{{\ensuremath{\Pi}}}1
    {Σ}{{\ensuremath{\Sigma}}}1
    {Φ}{{\ensuremath{\Phi}}}1
    {Ψ}{{\ensuremath{\Psi}}}1
    {Ω}{{\ensuremath{\Omega}}}1
    {ℵ}{{\ensuremath{\aleph}}}1
    {∀}{{\ensuremath{\forall}}}1
    {∃!}{{\ensuremath{\exists}!}}1
    {∃}{{\ensuremath{\exists}}}1
    {∈}{{\ensuremath{\in}}}1
    {∉}{{\ensuremath{\notin}}}1
    {∘}{{\ensuremath{\circ}}}1
    {∩}{{\ensuremath{\cap}}}1
    {∪}{{\ensuremath{\cup}}}1
    {⊆}{{\ensuremath{\subseteq}}}1
    {⊂}{{\ensuremath{\subset}}}1
    {⊇}{{\ensuremath{\supseteq}}}1
    {⊃}{{\ensuremath{\supset}}}1
    {⊄}{{\ensuremath{\nsubseteq}}}1
    {∅}{{\ensuremath{\emptyset}}}1
    {∧}{{\ensuremath{\wedge}}}1
    {∨}{{\ensuremath{\vee}}}1
    {¬}{{\ensuremath{\neg}}}1
    {⊤}{{\ensuremath{\top}}}1
    {⊥}{{\ensuremath{\bot}}}1
    {↔}{{\ensuremath{\leftrightarrow}}}1
    {→}{{\ensuremath{\rightarrow}}}1
    {←}{{\ensuremath{\leftarrow}}}1
    {↑}{{\ensuremath{\uparrow}}}1
    {⇒}{{\ensuremath{\Rightarrow}}}1
    {⇐}{{\ensuremath{\Leftarrow}}}1
    {⟹}{{\ensuremath{\Longrightarrow}}}1
    {↦}{{\ensuremath{\mapsto}}}1
    {≤}{{\ensuremath{\leq}}}1
    {≥}{{\ensuremath{\geq}}}1
    {≠}{{\ensuremath{\neq}}}1
    {≈}{{\ensuremath{\approx}}}1
    {≡}{{\ensuremath{\equiv}}}1
    {≃}{{\ensuremath{\simeq}}}1
    {≅}{{\ensuremath{\cong}}}1
    {×}{{\ensuremath{\times}}}1
    {·}{{\ensuremath{\cdot}}}1
    {±}{{\ensuremath{\pm}}}1
    {⊕}{{\ensuremath{\oplus}}}1
    {⊗}{{\ensuremath{\otimes}}}1
    {⋆}{{\ensuremath{\star}}}1
    {▸}{{\ensuremath{\blacktriangleright}}}1
    {⟨}{{\ensuremath{\langle}}}1
    {⟩}{{\ensuremath{\rangle}}}1
    {⦃}{{\ensuremath{\{\!\mid}}}1
    {⦄}{{\ensuremath{\mid\!\}}}}1
    {⟦}{{\ensuremath{[\![}}}1
    {⟧}{{\ensuremath{]\!]}}}1
    {⌊}{{\ensuremath{\lfloor}}}1
    {⌋}{{\ensuremath{\rfloor}}}1
    {⌈}{{\ensuremath{\lceil}}}1
    {⌉}{{\ensuremath{\rceil}}}1
    {√}{{\ensuremath{\surd}}}1
    {∣}{{\ensuremath{\mid}}}1
    {∤}{{\ensuremath{\nmid}}}1
    {∎}{{\ensuremath{\blacksquare}}}1
    {⋯}{{\ensuremath{\cdots}}}1
    {…}{{\ensuremath{\ldots}}}1
    {∑}{{\ensuremath{\sum}}}1
    {∏}{{\ensuremath{\prod}}}1
    {⁻¹}{{\ensuremath{^{-1}}}}1
    {¹}{{\ensuremath{^1}}}1
    {²}{{\ensuremath{^2}}}1
    {³}{{\ensuremath{^3}}}1
    {ⁿ}{{\ensuremath{^n}}}1
    {ᵐ}{{\ensuremath{^m}}}1
    {₀}{{\ensuremath{_0}}}1
    {₁}{{\ensuremath{_1}}}1
    {₂}{{\ensuremath{_2}}}1
    {₃}{{\ensuremath{_3}}}1
    {₄}{{\ensuremath{_4}}}1
    {₅}{{\ensuremath{_5}}}1
    {ₙ}{{\ensuremath{_n}}}1
    {ₘ}{{\ensuremath{_m}}}1
    {ₖ}{{\ensuremath{_k}}}1
    {ᵢ}{{\ensuremath{_i}}}1
    {ⱼ}{{\ensuremath{_j}}}1
    {ℤ}{{\ensuremath{\mathbb{Z}}}}1
    {ℕ}{{\ensuremath{\mathbb{N}}}}1
    {ℚ}{{\ensuremath{\mathbb{Q}}}}1
    {ℝ}{{\ensuremath{\mathbb{R}}}}1
    {ℂ}{{\ensuremath{\mathbb{C}}}}1
    {𝔽}{{\ensuremath{\mathbb{F}}}}1
    {𝒫}{{\ensuremath{\mathcal{P}}}}1
    {∗}{{\ensuremath{\ast}}}1
    {•}{{\ensuremath{\bullet}}}1,
  % Comments
  morecomment=[l]{--},
  morecomment=[s]{/-}{-/},
  morestring=[b]",
  stringstyle=\color{stringcolor},
  keepspaces=true,
  keywordstyle=[1]\color{keywordcolor}\bfseries,
  keywordstyle=[2]\color{sortcolor}\bfseries,
  keywordstyle=[3]\color{tacticcolor}\bfseries,
  keywordstyle=[4]\color{errorcolor}\bfseries\underline,
  escapeinside={(*@}{@*)},
  columns=fullflexible,
  basicstyle=\ttfamily\small,
  commentstyle=\color{commentcolor}\itshape,
  breaklines=true,
  showstringspaces=false,
  frame=single,
  framesep=2mm,
  xleftmargin=4mm,
  numbers=left,
  numberstyle=\tiny\color{commentcolor},
  numbersep=6pt,
}

% Inline Lean code in prose: \lean{Int.ModEq.pow}
\newcommand{\lean}[1]{\lstinline[language=lean,basicstyle=\ttfamily\normalsize]|#1|}

\lstset{language=lean}

% --- local rendering patch (2026-09-29) ---
% tectonic/XeTeX does not apply the listings literate table to multi-byte
% characters, so the Unicode set used by the listings of THIS note is mapped
% to math glyphs as active characters via the standard \lccode/\lowercase
% idiom (ASCII-only source, no extra packages). Characters deliberately NOT
% mapped: U+00B7 (middle dot) and U+2014 (em-dash), which xeCJK sets up at
% load time (making them active first aborts the run with "Control sequence
% ... already defined"; observed 2026-09-27 in the pendant note); they route
% through xeCJK + SimSun, as do all CJK ideographs and CJK punctuation.
% New relative to the template set (they occur in this note's sources):
% 00A7 (section sign), 00AB/00BB (guillemets, case labels), 21A6 (mapsto),
% 2208 (in), 2212 (minus sign), 27FA (long left-right arrow, iff).
% Rendering only: listing source bytes are untouched, so the byte-exactness
% checks are unaffected, and the \ifdefined guard leaves the pdflatex path
% identical to the template.
\ifdefined\XeTeXversion
  {\lccode`\~="00A7 \lowercase{\global\catcode`~=\active \global\def~{\S}}}
  {\lccode`\~="00AB \lowercase{\global\catcode`~=\active \global\def~{\guillemotleft}}}
  {\lccode`\~="00BB \lowercase{\global\catcode`~=\active \global\def~{\guillemotright}}}
  {\lccode`\~="2115 \lowercase{\global\catcode`~=\active \global\def~{\ensuremath{\mathbb{N}}}}}
  {\lccode`\~="2124 \lowercase{\global\catcode`~=\active \global\def~{\ensuremath{\mathbb{Z}}}}}
  {\lccode`\~="2192 \lowercase{\global\catcode`~=\active \global\def~{\ensuremath{\rightarrow}}}}
  {\lccode`\~="21A6 \lowercase{\global\catcode`~=\active \global\def~{\ensuremath{\mapsto}}}}
  {\lccode`\~="2200 \lowercase{\global\catcode`~=\active \global\def~{\ensuremath{\forall}}}}
  {\lccode`\~="2208 \lowercase{\global\catcode`~=\active \global\def~{\ensuremath{\in}}}}
  {\lccode`\~="2212 \lowercase{\global\catcode`~=\active \global\def~{\ensuremath{-}}}}
  {\lccode`\~="2227 \lowercase{\global\catcode`~=\active \global\def~{\ensuremath{\wedge}}}}
  {\lccode`\~="2265 \lowercase{\global\catcode`~=\active \global\def~{\ensuremath{\geq}}}}
  {\lccode`\~="27E8 \lowercase{\global\catcode`~=\active \global\def~{\ensuremath{\langle}}}}
  {\lccode`\~="27E9 \lowercase{\global\catcode`~=\active \global\def~{\ensuremath{\rangle}}}}
  {\lccode`\~="27FA \lowercase{\global\catcode`~=\active \global\def~{\ensuremath{\Longleftrightarrow}}}}
\fi
\ifdefined\XeTeXversion
  \usepackage{xeCJK}
  \setCJKmainfont{SimSun}
\fi

\newtheorem{theorem}{Theorem}

\title{Matchings of the period-3 pendant ladder: a mod-3 interlacing
recurrence and explicit modular residue tables\\ (Lean 4 machine-checked
claim of record)}
\author{Aurora0134\thanks{Contact: via the GitHub account Aurora0134; ORCID
placeholder --- to be filled in by the author before any upload. No personal
information is carried in this note.}}
\date{September 29, 2026}

\begin{document}
\maketitle

\begin{abstract}
Let $a(n)$ be the number of matchings (the empty one included) of the
period-3 pendant ladder: the ladder graph $P_n\square K_2$ with one pendant
leaf attached to the upper-rail vertex of every column $3k+1$. The AI4Math
pipeline formalised and machine-proved three exact laws of the
recurrence-defined counterpart $\mathtt{ap3}$ of $a$ (with
$\mathtt{ap3}\,n=a(n{+}1)$): (1)~a mod-3 interlacing identity that splits the
order-12 sparse recurrence into three order-4 subsequences with the same
coefficients $(45,63,11,-1)$; (2)~an explicit period-15 residue table for
$\mathtt{ap3}\bmod 2$; (3)~an explicit period-30 residue table for
$\mathtt{ap3}\bmod 4$. All three theorems compile with no \texttt{sorry},
and \texttt{\#print axioms} reports only \{propext, Quot.sound,
Classical.choice\}. Two independent novelty reviews found no occupying
record for the sequence, its recurrence, or its three mod-3 subsequences;
the plain ladder, the period-1 pendant (corona) and the period-2 pendant
relatives are registered, period 3 is not. The claim is strictly
\emph{new sequence / new recurrence}, never new mathematics. This note is a
priority claim of record, not a full paper.
\end{abstract}

\section{Introduction}

Three theorems about one integer sequence are reported here. The sequence
counts all matchings of the period-3 pendant ladder; its first values are
$a(0)=1$ (the empty graph with its empty matching) and, for $n\ge 1$,
$3,10,32,150,457,1489,6948,21191,69032,322142,\ldots$, and the conjectured
law recorded with the object was the sparse recurrence
$a(n)=45a(n-3)+63a(n-6)+11a(n-9)-a(n-12)$, equivalently: the three
$\bmod\,3$ subsequences all obey the same order-4 recurrence with
coefficients $(45,63,11,-1)$. The proofs were produced by the AI4Math
pipeline, in which LLM workers generate statements and proof scripts and the
Lean 4 kernel is the sole adjudicator at every gate: statements are frozen
by hash before proving, and a separate audit re-computes the symmetric
difference, scans the whole file for \texttt{sorry}/\texttt{admit} including
comments, and re-prints the axiom dependencies \cite{lean4}. The kernel
layer proves the pure-sequence statements (definitions below are
initial-values-plus-recurrence, $\mathbb{N}\to\mathbb{Z}$, no graph objects);
the identification of the recurrence-defined sequence with the matching
count of the pendant ladder is a conjecture layer carried by the pipeline's
computational evidence (five methodologically distinct counting
implementations in agreement over 241 terms, double-window modular scans)
and is not kernel-proved; Section~\ref{sec:limits} states this boundary
verbatim. This note is a priority claim of record deposited so that a dated,
citable public record exists ahead of the narrative paper; no literature
review is attempted here.

\section{Statement}\label{sec:statement}

Throughout, \texttt{ap3} is the main column: an
$\mathbb{N}\to\mathbb{Z}$ sequence with twelve frozen initial values and the
order-12 sparse recurrence
$\mathtt{ap3}\,(n+12)=45\,\mathtt{ap3}\,(n+9)+63\,\mathtt{ap3}\,(n+6)
+11\,\mathtt{ap3}\,(n+3)-\mathtt{ap3}\,n$. The index convention is fixed by
the pipeline's claims record: $\mathtt{ap3}\,n=a(n{+}1)$ (0-based Lean
against 1-based counting), the twelve frozen initial values are the 0-based
$a(0..11)$ (the base term $a(0)=1$ included), and \texttt{bp0}, \texttt{bp1},
\texttt{bp2} are the three $\bmod\,3$ subsequences, $\mathtt{bp}_r\,k
=a(3k{+}r{+}1)$ in 1-based terms, each with four initial values (the
$\mathtt{ap3}$ values at indices $3k+r$) and the order-4 recurrence
$\mathtt{bp}_r\,(k+4)=45\,\mathtt{bp}_r\,(k+3)+63\,\mathtt{bp}_r\,(k+2)
+11\,\mathtt{bp}_r\,(k+1)-\mathtt{bp}_r\,k$. \texttt{pat15} and
\texttt{pat30} are the explicit residue tables: $\mathbb{N}\to\mathbb{Z}$
functions equal to $0$ outside the listed classes.

% LEAN: tasks/20260929-comb10-ladder-pend3 :: ap3_interleave grade=完全证明
\begin{theorem}[mod-3 interlacing identity]\label{thm:interleave}
For every $k:\mathbb{N}$,
\[
\mathtt{ap3}\,(3k)=\mathtt{bp0}\,k
\;\wedge\;\mathtt{ap3}\,(3k+1)=\mathtt{bp1}\,k
\;\wedge\;\mathtt{ap3}\,(3k+2)=\mathtt{bp2}\,k .
\]
That is: the order-12 sparse recurrence decouples along $n\bmod 3$ into
three order-4 subsequences that share the coefficient vector
$(45,63,11,-1)$.
\end{theorem}

% LEAN: tasks/20260929-comb10-ladder-pend3 :: ap3_mod2_period15 grade=完全证明
\begin{theorem}[explicit period-15 residue table for mod 2]\label{thm:mod2}
For all $n,r:\mathbb{N}$ with $n \bmod 15 = r$,
$\mathtt{ap3}\,n \bmod 2 = \mathtt{pat15}\,r$, where \texttt{pat15} is the
explicit table $(1,1,0,0,0,1,1,0,1,0,0,0,0,1,0)$ on the fifteen residue
classes; equivalently $\mathtt{ap3}\,n$ is odd exactly when
$n \bmod 15 \in \{0,1,5,6,8,13\}$.
\end{theorem}

% LEAN: tasks/20260929-comb10-ladder-pend3 :: ap3_mod4_period30 grade=完全证明
\begin{theorem}[explicit period-30 residue table for mod 4]\label{thm:mod4}
For all $n,r:\mathbb{N}$ with $n \bmod 30 = r$,
$\mathtt{ap3}\,n \bmod 4 = \mathtt{pat30}\,r$, where \texttt{pat30} is the
explicit thirty-entry table reproduced in the listing below.
\end{theorem}

\begin{lstlisting}[caption={Frozen base definitions and the head of
Theorem~\ref{thm:interleave}, verbatim lines 37--89 of
\texttt{formalized/comb10-ladder-pend3.statement.txt} (sha256
\texttt{47d54db7}): the four sequence definitions and the interlacing
statement.},label={lst:s1}]
/-- **序列 ap3（主列，12 阶稀疏递推）**。
初值（12 项，0-based a(0..11)）：1, 3, 10, 32, 150, 457, 1489, 6948, 21191, 69032, 322142, 982496；
递推 ap3 (n+12) = 45·ap3 (n+9) + 63·ap3 (n+6) + 11·ap3 (n+3) − ap3 n（n ≥ 0）。 -/
def ap3 : ℕ → ℤ
  | 0 => 1
  | 1 => 3
  | 2 => 10
  | 3 => 32
  | 4 => 150
  | 5 => 457
  | 6 => 1489
  | 7 => 6948
  | 8 => 21191
  | 9 => 69032
  | 10 => 322142
  | 11 => 982496
  | n + 12 => 45 * ap3 (n + 9) + 63 * ap3 (n + 6) + 11 * ap3 (n + 3) - ap3 n

/-- **序列 bp0（mod-3 子列 r = 0，四阶递推）**。
初值 = ap3 在下标 0, 3, 6, 9 处之值：1, 32, 1489, 69032；
递推 bp0 (k+4) = 45·bp0 (k+3) + 63·bp0 (k+2) + 11·bp0 (k+1) − bp0 k（k ≥ 0）。 -/
def bp0 : ℕ → ℤ
  | 0 => 1
  | 1 => 32
  | 2 => 1489
  | 3 => 69032
  | k + 4 => 45 * bp0 (k + 3) + 63 * bp0 (k + 2) + 11 * bp0 (k + 1) - bp0 k

/-- **序列 bp1（mod-3 子列 r = 1，四阶递推）**。
初值 = ap3 在下标 1, 4, 7, 10 处之值：3, 150, 6948, 322142；
递推 bp1 (k+4) = 45·bp1 (k+3) + 63·bp1 (k+2) + 11·bp1 (k+1) − bp1 k（k ≥ 0）。 -/
def bp1 : ℕ → ℤ
  | 0 => 3
  | 1 => 150
  | 2 => 6948
  | 3 => 322142
  | k + 4 => 45 * bp1 (k + 3) + 63 * bp1 (k + 2) + 11 * bp1 (k + 1) - bp1 k

/-- **序列 bp2（mod-3 子列 r = 2，四阶递推）**。
初值 = ap3 在下标 2, 5, 8, 11 处之值：10, 457, 21191, 982496；
递推 bp2 (k+4) = 45·bp2 (k+3) + 63·bp2 (k+2) + 11·bp2 (k+1) − bp2 k（k ≥ 0）。 -/
def bp2 : ℕ → ℤ
  | 0 => 10
  | 1 => 457
  | 2 => 21191
  | 3 => 982496
  | k + 4 => 45 * bp2 (k + 3) + 63 * bp2 (k + 2) + 11 * bp2 (k + 1) - bp2 k

/-- **T1（主攻）**：主列 12 阶稀疏递推 = 三条同系数四阶子列按 n mod 3 交织——
对一切 k，ap3 (3k) = bp0 k、ap3 (3k+1) = bp1 k、ap3 (3k+2) = bp2 k。
（claims.md §T1：递推只含 3 的倍数滞后，三条 mod-3 副列天然解耦，无需辅助 12 阶恒等式引理。） -/
theorem ap3_interleave (k : ℕ) :
    ap3 (3 * k) = bp0 k ∧ ap3 (3 * k + 1) = bp1 k ∧ ap3 (3 * k + 2) = bp2 k := by
\end{lstlisting}

\begin{lstlisting}[caption={Frozen definition of \texttt{pat15} and the head
of Theorem~\ref{thm:mod2}, verbatim lines 92--116 of the frozen statement
file.},label={lst:s2}]
/-- **余数样式 pat15**（15 留数类 ↦ mod 2 余数）：
0↦1, 1↦1, 2↦0, 3↦0, 4↦0, 5↦1, 6↦1, 7↦0, 8↦1, 9↦0,
10↦0, 11↦0, 12↦0, 13↦1，其余 ↦ 0
（即 ap3 n 为奇 ⟺ n mod 15 ∈ {0, 1, 5, 6, 8, 13}）。 -/
def pat15 : ℕ → ℤ
  | 0 => 1
  | 1 => 1
  | 2 => 0
  | 3 => 0
  | 4 => 0
  | 5 => 1
  | 6 => 1
  | 7 => 0
  | 8 => 1
  | 9 => 0
  | 10 => 0
  | 11 => 0
  | 12 => 0
  | 13 => 1
  | _ => 0

/-- **T2（伴随）**：ap3 的 mod 2 余数仅依赖 n mod 15，样式为 pat15
（0-based 口径；claims.md §T2 样式 (1,1,0,0,0,1,1,0,1,0,0,0,0,1,0)）。 -/
theorem ap3_mod2_period15 (n r : ℕ) (hr : n % 15 = r) :
    ap3 n % 2 = pat15 r := by
\end{lstlisting}

\begin{lstlisting}[caption={Frozen definition of \texttt{pat30} and the head
of Theorem~\ref{thm:mod4}, verbatim lines 119--159 of the frozen statement
file.},label={lst:s3}]
/-- **余数样式 pat30**（30 留数类 ↦ mod 4 余数）：
0↦1, 1↦3, 2↦2, 3↦0, 4↦2, 5↦1, 6↦1, 7↦0, 8↦3, 9↦0,
10↦2, 11↦0, 12↦2, 13↦1, 14↦2, 15↦1, 16↦1, 17↦2, 18↦2, 19↦2,
20↦1, 21↦3, 22↦2, 23↦1, 24↦2, 25↦2, 26↦0, 27↦0, 28↦1，其余 ↦ 0。 -/
def pat30 : ℕ → ℤ
  | 0 => 1
  | 1 => 3
  | 2 => 2
  | 3 => 0
  | 4 => 2
  | 5 => 1
  | 6 => 1
  | 7 => 0
  | 8 => 3
  | 9 => 0
  | 10 => 2
  | 11 => 0
  | 12 => 2
  | 13 => 1
  | 14 => 2
  | 15 => 1
  | 16 => 1
  | 17 => 2
  | 18 => 2
  | 19 => 2
  | 20 => 1
  | 21 => 3
  | 22 => 2
  | 23 => 1
  | 24 => 2
  | 25 => 2
  | 26 => 0
  | 27 => 0
  | 28 => 1
  | _ => 0

/-- **T3（伴随）**：ap3 的 mod 4 余数仅依赖 n mod 30，样式为 pat30
（蕴含 T2：奇值类 mod 30 = {0, 1, 5, 6, 8, 13, 15, 16, 20, 21, 23, 28}，每类 mod 2 与 pat15 一致；
蕴含关系归证明侧，statement 层 T2/T3 并列独立）。 -/
theorem ap3_mod4_period30 (n r : ℕ) (hr : n % 30 = r) :
    ap3 n % 4 = pat30 r := by
\end{lstlisting}

The complete final file is reproduced verbatim in Section~\ref{sec:proof}
(whole file, no excerpting); it also ships in this deposit under
\texttt{proofs/}. The frozen statement file carries the three theorem
bodies as \texttt{sorry} placeholders (that is its role: the formalisation
department freezes statements, the proving department replaces bodies); the
excerpts above stop at each theorem head, so no \texttt{sorry} line enters
any listing, and the final proof file has none anywhere.

\section{Proof}\label{sec:proof}

All three proofs are strong inductions on the index, with the recursion
depth 12 (the sparse recurrence's gap) between the induction hypothesis
instances used and the conclusion. For Theorem~\ref{thm:interleave} the
decoupling is structural: the sparse recurrence only lags by multiples of
$3$, so each $\bmod\,3$ subsequence satisfies the order-4 recurrence on its
own; the proof first establishes the three one-sided lemmas
$\mathtt{ap3}\,(3k+r)=\mathtt{bp}_r\,k$ ($r=0,1,2$) by strong induction
(twelve base cases discharged by case enumeration and arithmetic normalisation;
induction step: unfold $\mathtt{bp}_r$, rewrite the four occurrences through
the induction hypotheses, and match the $\mathtt{ap3}$ recurrence instance at
index $3k+r$), then assembles the conjunction. Theorems~\ref{thm:mod2}
and~\ref{thm:mod4} prove the full-column residue laws directly: a
$\mathtt{period}$ lemma ($\forall m$, $\mathtt{ap3}\,m \bmod q =
\mathtt{pat}\,(m \bmod P)$) is established by strong induction, where the
induction step rewrites the recurrence modulo $q$ (for $q=2$: all four
coefficients are $1$; for $q=4$: $1,3,3,3$), enumerates the $P$ residue
classes of the induction target $m \bmod P$ ($P=15$, resp.\ $30$), shifts
each of the four lagged terms into a decided residue class, and closes each
case by table lookup and arithmetic; the outer theorem instantiates the
$\mathtt{period}$ lemma at $m=n$. The third theorem's file carries a
\texttt{set\_option maxHeartbeats 800000} decoration: an elaboration-resource
knob for the thirty-case enumeration, with no effect on the statement or the
kernel semantics. The complete final file follows verbatim ($440$ lines, no
excerpting; the working copy is CRLF, listing LF-normalised per the
chorded-cycle precedent, byte-exact original in \texttt{proofs/}).

\begin{lstlisting}[basicstyle=\ttfamily\footnotesize,caption={Complete proof
file of Theorems~\ref{thm:interleave}, \ref{thm:mod2} and \ref{thm:mod4},
verbatim lines 1--440 (whole file) of
\texttt{final/comb10-ladder-pend3-proved.lean} (sha256 \texttt{e56e496f}).},label={lst:proof}]
/-
  AI4Math 流水线 · AI 生成 · 2026-09-29
  任务线：comb10-ladder-pend3（pool-comb-10）· 三定理合并终审稿
  冻结快照：formalized/comb10-ladder-pend3.statement.txt
    （sha256 47d54db7c3d5309720537295055d5869e990618399950aaacb311f8028a9bba9）
  来源分件（kernel 已逐件过）：attempts/comb10-t1.lean / comb10-t2.lean / comb10-t3.lean
  声明层对冻结件零改动；T3 前一条 set_option maxHeartbeats 800000 为编译资源旋钮
  （comb-09 T2 同类判例，内核语义不受影响）。
-/

import Mathlib

def ap3 : ℕ → ℤ
  | 0 => 1
  | 1 => 3
  | 2 => 10
  | 3 => 32
  | 4 => 150
  | 5 => 457
  | 6 => 1489
  | 7 => 6948
  | 8 => 21191
  | 9 => 69032
  | 10 => 322142
  | 11 => 982496
  | n + 12 => 45 * ap3 (n + 9) + 63 * ap3 (n + 6) + 11 * ap3 (n + 3) - ap3 n

def bp0 : ℕ → ℤ
  | 0 => 1
  | 1 => 32
  | 2 => 1489
  | 3 => 69032
  | k + 4 => 45 * bp0 (k + 3) + 63 * bp0 (k + 2) + 11 * bp0 (k + 1) - bp0 k

def bp1 : ℕ → ℤ
  | 0 => 3
  | 1 => 150
  | 2 => 6948
  | 3 => 322142
  | k + 4 => 45 * bp1 (k + 3) + 63 * bp1 (k + 2) + 11 * bp1 (k + 1) - bp1 k

def bp2 : ℕ → ℤ
  | 0 => 10
  | 1 => 457
  | 2 => 21191
  | 3 => 982496
  | k + 4 => 45 * bp2 (k + 3) + 63 * bp2 (k + 2) + 11 * bp2 (k + 1) - bp2 k

theorem ap3_sub_0 (k : ℕ) : ap3 (3 * k) = bp0 k := by
  induction k using Nat.strong_induction_on with
  | h k ih =>
    by_cases hk : k < 4
    · interval_cases k <;> norm_num [ap3, bp0]
    · let t := k - 4
      have ht : k = t + 4 := by dsimp [t]; omega
      have h0 := ih t (by omega)
      have h1 := ih (t + 1) (by omega)
      have h2 := ih (t + 2) (by omega)
      have h3 := ih (t + 3) (by omega)
      rw [ht, bp0]
      have hidx : 3 * (t + 4) = 3 * t + 12 := by omega
      rw [hidx]
      have hrec : ap3 (3 * t + 12) =
          45 * ap3 (3 * t + 9) + 63 * ap3 (3 * t + 6) +
            11 * ap3 (3 * t + 3) - ap3 (3 * t) := by
        rw [ap3]
      rw [hrec]
      have e3 : 3 * t + 9 = 3 * (t + 3) := by omega
      have e2 : 3 * t + 6 = 3 * (t + 2) := by omega
      have e1 : 3 * t + 3 = 3 * (t + 1) := by omega
      rw [e3, e2, e1, h3, h2, h1, h0]

theorem ap3_sub_1 (k : ℕ) : ap3 (3 * k + 1) = bp1 k := by
  induction k using Nat.strong_induction_on with
  | h k ih =>
    by_cases hk : k < 4
    · interval_cases k <;> norm_num [ap3, bp1]
    · let t := k - 4
      have ht : k = t + 4 := by dsimp [t]; omega
      have h0 := ih t (by omega)
      have h1 := ih (t + 1) (by omega)
      have h2 := ih (t + 2) (by omega)
      have h3 := ih (t + 3) (by omega)
      rw [ht, bp1]
      have hidx : 3 * (t + 4) + 1 = (3 * t + 1) + 12 := by omega
      rw [hidx]
      have hrec : ap3 ((3 * t + 1) + 12) =
          45 * ap3 ((3 * t + 1) + 9) + 63 * ap3 ((3 * t + 1) + 6) +
            11 * ap3 ((3 * t + 1) + 3) - ap3 (3 * t + 1) := by
        rw [ap3]
      rw [hrec]
      have e3 : (3 * t + 1) + 9 = 3 * (t + 3) + 1 := by omega
      have e2 : (3 * t + 1) + 6 = 3 * (t + 2) + 1 := by omega
      have e1 : (3 * t + 1) + 3 = 3 * (t + 1) + 1 := by omega
      rw [e3, e2, e1, h3, h2, h1, h0]

theorem ap3_sub_2 (k : ℕ) : ap3 (3 * k + 2) = bp2 k := by
  induction k using Nat.strong_induction_on with
  | h k ih =>
    by_cases hk : k < 4
    · interval_cases k <;> norm_num [ap3, bp2]
    · let t := k - 4
      have ht : k = t + 4 := by dsimp [t]; omega
      have h0 := ih t (by omega)
      have h1 := ih (t + 1) (by omega)
      have h2 := ih (t + 2) (by omega)
      have h3 := ih (t + 3) (by omega)
      rw [ht, bp2]
      have hidx : 3 * (t + 4) + 2 = (3 * t + 2) + 12 := by omega
      rw [hidx]
      have hrec : ap3 ((3 * t + 2) + 12) =
          45 * ap3 ((3 * t + 2) + 9) + 63 * ap3 ((3 * t + 2) + 6) +
            11 * ap3 ((3 * t + 2) + 3) - ap3 (3 * t + 2) := by
        rw [ap3]
      rw [hrec]
      have e3 : (3 * t + 2) + 9 = 3 * (t + 3) + 2 := by omega
      have e2 : (3 * t + 2) + 6 = 3 * (t + 2) + 2 := by omega
      have e1 : (3 * t + 2) + 3 = 3 * (t + 1) + 2 := by omega
      rw [e3, e2, e1, h3, h2, h1, h0]

theorem ap3_interleave (k : ℕ) :
    ap3 (3 * k) = bp0 k ∧ ap3 (3 * k + 1) = bp1 k ∧ ap3 (3 * k + 2) = bp2 k := by
  exact ⟨ap3_sub_0 k, ap3_sub_1 k, ap3_sub_2 k⟩

def pat15 : ℕ → ℤ
  | 0 => 1
  | 1 => 1
  | 2 => 0
  | 3 => 0
  | 4 => 0
  | 5 => 1
  | 6 => 1
  | 7 => 0
  | 8 => 1
  | 9 => 0
  | 10 => 0
  | 11 => 0
  | 12 => 0
  | 13 => 1
  | _ => 0

theorem ap3_mod2_period15 (n r : ℕ) (hr : n % 15 = r) :
    ap3 n % 2 = pat15 r := by
  have hperiod : ∀ m : ℕ, ap3 m % 2 = pat15 (m % 15) := by
    intro m
    induction m using Nat.strong_induction_on with
    | h m ih =>
      by_cases hm : m < 12
      · interval_cases m <;> decide
      · let t := m - 12
        have hmt : m = t + 12 := by
          dsimp [t]
          omega
        rw [hmt, ap3]
        have h0 := ih t (by omega)
        have h3 := ih (t + 3) (by omega)
        have h6 := ih (t + 6) (by omega)
        have h9 := ih (t + 9) (by omega)
        have hresBound : t % 15 < 15 := Nat.mod_lt _ (by omega)
        interval_cases hres : t % 15
        case «0» =>
          rw [(show (t + 3) % 15 = 3 by omega), (show (t + 6) % 15 = 6 by omega),
              (show (t + 9) % 15 = 9 by omega), (show (t + 12) % 15 = 12 by omega)] at *
          norm_num [pat15] at *
          omega
        case «1» =>
          rw [(show (t + 3) % 15 = 4 by omega), (show (t + 6) % 15 = 7 by omega),
              (show (t + 9) % 15 = 10 by omega), (show (t + 12) % 15 = 13 by omega)] at *
          norm_num [pat15] at *
          omega
        case «2» =>
          rw [(show (t + 3) % 15 = 5 by omega), (show (t + 6) % 15 = 8 by omega),
              (show (t + 9) % 15 = 11 by omega), (show (t + 12) % 15 = 14 by omega)] at *
          norm_num [pat15] at *
          omega
        case «3» =>
          rw [(show (t + 3) % 15 = 6 by omega), (show (t + 6) % 15 = 9 by omega),
              (show (t + 9) % 15 = 12 by omega), (show (t + 12) % 15 = 0 by omega)] at *
          norm_num [pat15] at *
          omega
        case «4» =>
          rw [(show (t + 3) % 15 = 7 by omega), (show (t + 6) % 15 = 10 by omega),
              (show (t + 9) % 15 = 13 by omega), (show (t + 12) % 15 = 1 by omega)] at *
          norm_num [pat15] at *
          omega
        case «5» =>
          rw [(show (t + 3) % 15 = 8 by omega), (show (t + 6) % 15 = 11 by omega),
              (show (t + 9) % 15 = 14 by omega), (show (t + 12) % 15 = 2 by omega)] at *
          norm_num [pat15] at *
          omega
        case «6» =>
          rw [(show (t + 3) % 15 = 9 by omega), (show (t + 6) % 15 = 12 by omega),
              (show (t + 9) % 15 = 0 by omega), (show (t + 12) % 15 = 3 by omega)] at *
          norm_num [pat15] at *
          omega
        case «7» =>
          rw [(show (t + 3) % 15 = 10 by omega), (show (t + 6) % 15 = 13 by omega),
              (show (t + 9) % 15 = 1 by omega), (show (t + 12) % 15 = 4 by omega)] at *
          norm_num [pat15] at *
          omega
        case «8» =>
          rw [(show (t + 3) % 15 = 11 by omega), (show (t + 6) % 15 = 14 by omega),
              (show (t + 9) % 15 = 2 by omega), (show (t + 12) % 15 = 5 by omega)] at *
          norm_num [pat15] at *
          omega
        case «9» =>
          rw [(show (t + 3) % 15 = 12 by omega), (show (t + 6) % 15 = 0 by omega),
              (show (t + 9) % 15 = 3 by omega), (show (t + 12) % 15 = 6 by omega)] at *
          norm_num [pat15] at *
          omega
        case «10» =>
          rw [(show (t + 3) % 15 = 13 by omega), (show (t + 6) % 15 = 1 by omega),
              (show (t + 9) % 15 = 4 by omega), (show (t + 12) % 15 = 7 by omega)] at *
          norm_num [pat15] at *
          omega
        case «11» =>
          rw [(show (t + 3) % 15 = 14 by omega), (show (t + 6) % 15 = 2 by omega),
              (show (t + 9) % 15 = 5 by omega), (show (t + 12) % 15 = 8 by omega)] at *
          norm_num [pat15] at *
          omega
        case «12» =>
          rw [(show (t + 3) % 15 = 0 by omega), (show (t + 6) % 15 = 3 by omega),
              (show (t + 9) % 15 = 6 by omega), (show (t + 12) % 15 = 9 by omega)] at *
          norm_num [pat15] at *
          omega
        case «13» =>
          rw [(show (t + 3) % 15 = 1 by omega), (show (t + 6) % 15 = 4 by omega),
              (show (t + 9) % 15 = 7 by omega), (show (t + 12) % 15 = 10 by omega)] at *
          norm_num [pat15] at *
          omega
        case «14» =>
          rw [(show (t + 3) % 15 = 2 by omega), (show (t + 6) % 15 = 5 by omega),
              (show (t + 9) % 15 = 8 by omega), (show (t + 12) % 15 = 11 by omega)] at *
          norm_num [pat15] at *
          omega
  rw [hperiod n, hr]

def pat30 : ℕ → ℤ
  | 0 => 1
  | 1 => 3
  | 2 => 2
  | 3 => 0
  | 4 => 2
  | 5 => 1
  | 6 => 1
  | 7 => 0
  | 8 => 3
  | 9 => 0
  | 10 => 2
  | 11 => 0
  | 12 => 2
  | 13 => 1
  | 14 => 2
  | 15 => 1
  | 16 => 1
  | 17 => 2
  | 18 => 2
  | 19 => 2
  | 20 => 1
  | 21 => 3
  | 22 => 2
  | 23 => 1
  | 24 => 2
  | 25 => 2
  | 26 => 0
  | 27 => 0
  | 28 => 1
  | _ => 0

set_option maxHeartbeats 800000 in
theorem ap3_mod4_period30 (n r : ℕ) (hr : n % 30 = r) :
    ap3 n % 4 = pat30 r := by
  have hperiod : ∀ m : ℕ, ap3 m % 4 = pat30 (m % 30) := by
    intro m
    induction m using Nat.strong_induction_on with
    | h m ih =>
      by_cases hm : m < 12
      · interval_cases m <;> decide
      · let t := m - 12
        have hmt : m = t + 12 := by
          dsimp [t]
          omega
        rw [hmt, ap3]
        have h0 := ih t (by omega)
        have h3 := ih (t + 3) (by omega)
        have h6 := ih (t + 6) (by omega)
        have h9 := ih (t + 9) (by omega)
        have hresBound : t % 30 < 30 := Nat.mod_lt _ (by omega)
        interval_cases hres : t % 30
        case «0» =>
          rw [(show (t + 3) % 30 = 3 by omega), (show (t + 6) % 30 = 6 by omega),
              (show (t + 9) % 30 = 9 by omega), (show (t + 12) % 30 = 12 by omega)] at *
          norm_num [pat30] at *
          omega
        case «1» =>
          rw [(show (t + 3) % 30 = 4 by omega), (show (t + 6) % 30 = 7 by omega),
              (show (t + 9) % 30 = 10 by omega), (show (t + 12) % 30 = 13 by omega)] at *
          norm_num [pat30] at *
          omega
        case «2» =>
          rw [(show (t + 3) % 30 = 5 by omega), (show (t + 6) % 30 = 8 by omega),
              (show (t + 9) % 30 = 11 by omega), (show (t + 12) % 30 = 14 by omega)] at *
          norm_num [pat30] at *
          omega
        case «3» =>
          rw [(show (t + 3) % 30 = 6 by omega), (show (t + 6) % 30 = 9 by omega),
              (show (t + 9) % 30 = 12 by omega), (show (t + 12) % 30 = 15 by omega)] at *
          norm_num [pat30] at *
          omega
        case «4» =>
          rw [(show (t + 3) % 30 = 7 by omega), (show (t + 6) % 30 = 10 by omega),
              (show (t + 9) % 30 = 13 by omega), (show (t + 12) % 30 = 16 by omega)] at *
          norm_num [pat30] at *
          omega
        case «5» =>
          rw [(show (t + 3) % 30 = 8 by omega), (show (t + 6) % 30 = 11 by omega),
              (show (t + 9) % 30 = 14 by omega), (show (t + 12) % 30 = 17 by omega)] at *
          norm_num [pat30] at *
          omega
        case «6» =>
          rw [(show (t + 3) % 30 = 9 by omega), (show (t + 6) % 30 = 12 by omega),
              (show (t + 9) % 30 = 15 by omega), (show (t + 12) % 30 = 18 by omega)] at *
          norm_num [pat30] at *
          omega
        case «7» =>
          rw [(show (t + 3) % 30 = 10 by omega), (show (t + 6) % 30 = 13 by omega),
              (show (t + 9) % 30 = 16 by omega), (show (t + 12) % 30 = 19 by omega)] at *
          norm_num [pat30] at *
          omega
        case «8» =>
          rw [(show (t + 3) % 30 = 11 by omega), (show (t + 6) % 30 = 14 by omega),
              (show (t + 9) % 30 = 17 by omega), (show (t + 12) % 30 = 20 by omega)] at *
          norm_num [pat30] at *
          omega
        case «9» =>
          rw [(show (t + 3) % 30 = 12 by omega), (show (t + 6) % 30 = 15 by omega),
              (show (t + 9) % 30 = 18 by omega), (show (t + 12) % 30 = 21 by omega)] at *
          norm_num [pat30] at *
          omega
        case «10» =>
          rw [(show (t + 3) % 30 = 13 by omega), (show (t + 6) % 30 = 16 by omega),
              (show (t + 9) % 30 = 19 by omega), (show (t + 12) % 30 = 22 by omega)] at *
          norm_num [pat30] at *
          omega
        case «11» =>
          rw [(show (t + 3) % 30 = 14 by omega), (show (t + 6) % 30 = 17 by omega),
              (show (t + 9) % 30 = 20 by omega), (show (t + 12) % 30 = 23 by omega)] at *
          norm_num [pat30] at *
          omega
        case «12» =>
          rw [(show (t + 3) % 30 = 15 by omega), (show (t + 6) % 30 = 18 by omega),
              (show (t + 9) % 30 = 21 by omega), (show (t + 12) % 30 = 24 by omega)] at *
          norm_num [pat30] at *
          omega
        case «13» =>
          rw [(show (t + 3) % 30 = 16 by omega), (show (t + 6) % 30 = 19 by omega),
              (show (t + 9) % 30 = 22 by omega), (show (t + 12) % 30 = 25 by omega)] at *
          norm_num [pat30] at *
          omega
        case «14» =>
          rw [(show (t + 3) % 30 = 17 by omega), (show (t + 6) % 30 = 20 by omega),
              (show (t + 9) % 30 = 23 by omega), (show (t + 12) % 30 = 26 by omega)] at *
          norm_num [pat30] at *
          omega
        case «15» =>
          rw [(show (t + 3) % 30 = 18 by omega), (show (t + 6) % 30 = 21 by omega),
              (show (t + 9) % 30 = 24 by omega), (show (t + 12) % 30 = 27 by omega)] at *
          norm_num [pat30] at *
          omega
        case «16» =>
          rw [(show (t + 3) % 30 = 19 by omega), (show (t + 6) % 30 = 22 by omega),
              (show (t + 9) % 30 = 25 by omega), (show (t + 12) % 30 = 28 by omega)] at *
          norm_num [pat30] at *
          omega
        case «17» =>
          rw [(show (t + 3) % 30 = 20 by omega), (show (t + 6) % 30 = 23 by omega),
              (show (t + 9) % 30 = 26 by omega), (show (t + 12) % 30 = 29 by omega)] at *
          norm_num [pat30] at *
          omega
        case «18» =>
          rw [(show (t + 3) % 30 = 21 by omega), (show (t + 6) % 30 = 24 by omega),
              (show (t + 9) % 30 = 27 by omega), (show (t + 12) % 30 = 0 by omega)] at *
          norm_num [pat30] at *
          omega
        case «19» =>
          rw [(show (t + 3) % 30 = 22 by omega), (show (t + 6) % 30 = 25 by omega),
              (show (t + 9) % 30 = 28 by omega), (show (t + 12) % 30 = 1 by omega)] at *
          norm_num [pat30] at *
          omega
        case «20» =>
          rw [(show (t + 3) % 30 = 23 by omega), (show (t + 6) % 30 = 26 by omega),
              (show (t + 9) % 30 = 29 by omega), (show (t + 12) % 30 = 2 by omega)] at *
          norm_num [pat30] at *
          omega
        case «21» =>
          rw [(show (t + 3) % 30 = 24 by omega), (show (t + 6) % 30 = 27 by omega),
              (show (t + 9) % 30 = 0 by omega), (show (t + 12) % 30 = 3 by omega)] at *
          norm_num [pat30] at *
          omega
        case «22» =>
          rw [(show (t + 3) % 30 = 25 by omega), (show (t + 6) % 30 = 28 by omega),
              (show (t + 9) % 30 = 1 by omega), (show (t + 12) % 30 = 4 by omega)] at *
          norm_num [pat30] at *
          omega
        case «23» =>
          rw [(show (t + 3) % 30 = 26 by omega), (show (t + 6) % 30 = 29 by omega),
              (show (t + 9) % 30 = 2 by omega), (show (t + 12) % 30 = 5 by omega)] at *
          norm_num [pat30] at *
          omega
        case «24» =>
          rw [(show (t + 3) % 30 = 27 by omega), (show (t + 6) % 30 = 0 by omega),
              (show (t + 9) % 30 = 3 by omega), (show (t + 12) % 30 = 6 by omega)] at *
          norm_num [pat30] at *
          omega
        case «25» =>
          rw [(show (t + 3) % 30 = 28 by omega), (show (t + 6) % 30 = 1 by omega),
              (show (t + 9) % 30 = 4 by omega), (show (t + 12) % 30 = 7 by omega)] at *
          norm_num [pat30] at *
          omega
        case «26» =>
          rw [(show (t + 3) % 30 = 29 by omega), (show (t + 6) % 30 = 2 by omega),
              (show (t + 9) % 30 = 5 by omega), (show (t + 12) % 30 = 8 by omega)] at *
          norm_num [pat30] at *
          omega
        case «27» =>
          rw [(show (t + 3) % 30 = 0 by omega), (show (t + 6) % 30 = 3 by omega),
              (show (t + 9) % 30 = 6 by omega), (show (t + 12) % 30 = 9 by omega)] at *
          norm_num [pat30] at *
          omega
        case «28» =>
          rw [(show (t + 3) % 30 = 1 by omega), (show (t + 6) % 30 = 4 by omega),
              (show (t + 9) % 30 = 7 by omega), (show (t + 12) % 30 = 10 by omega)] at *
          norm_num [pat30] at *
          omega
        case «29» =>
          rw [(show (t + 3) % 30 = 2 by omega), (show (t + 6) % 30 = 5 by omega),
              (show (t + 9) % 30 = 8 by omega), (show (t + 12) % 30 = 11 by omega)] at *
          norm_num [pat30] at *
          omega
  rw [hperiod n, hr]
\end{lstlisting}

\section{Verification evidence}

\begin{itemize}
\item Toolchain: Lean \texttt{v4.34.0} (\texttt{leanprover/lean4:v4.34.0}),
mathlib4 \texttt{v4.34.0}, revision \texttt{5ed29652}. Reproduction: with
that pin, \texttt{lake env lean} on
\texttt{proofs/01-comb10-ladder-pend3-proved.lean} compiles with exit code
$0$ and empty output.
\item Statement integrity: the frozen statement file (160 lines) has sha256
\begin{center}
\texttt{47d54db7c3d5309720537295055d5869e990618}\allowbreak
\texttt{399950aaacb311f8028a9bba9},
\end{center}
and the final proof file (440 lines) has sha256
\begin{center}
\texttt{e56e496f8ddc8310316a1d6f1f9befa17e76200}\allowbreak
\texttt{f97b22515dcd0b2260e9fd7b1}.
\end{center}
A re-run of the audit department's symmetric-difference method (comments
stripped, declarations split, definitions compared full-block, theorems up
to \texttt{:=}) reports all nine frozen declarations header-identical, with
exactly the three auxiliary lemmas \texttt{ap3\_sub\_0/1/2} added as
proof-body scaffolding (\texttt{audit/statement-diff.out}).
\item Kernel check: compiles with $0$ \texttt{sorry} (whole-file scan
including comments, also covering \texttt{admit}: zero hits);
\texttt{\#print axioms} output, verbatim (kernel re-run of 2026-09-29,
exit code 0, kept in \texttt{audit/axioms.out}):
\begin{lstlisting}[basicstyle=\ttfamily\footnotesize,frame=none]
'ap3_interleave' depends on axioms: [propext, Quot.sound]
'ap3_mod2_period15' depends on axioms: [propext, Classical.choice, Quot.sound]
'ap3_mod4_period30' depends on axioms: [propext, Classical.choice, Quot.sound]
\end{lstlisting}
Each line is a subset of the whitelist \{propext, Quot.sound,
Classical.choice\}.
\item Grading by the audit department: all three proven declarations
\emph{complete proof}; the lowest grade reached for the case = \emph{complete
proof} (final adjudication \texttt{final-comb10-ladder-pend3.txt}: statement
$9/9$ byte-identical, strict compile exit 0, zero \texttt{sorry}/%
\texttt{admit}, axiom whitelist subset). No weakening, no added hypotheses,
no unproven residue.
\item Audit chain (originals in this deposit under \texttt{audit/} unless
noted): the final adjudication, the statement-difference re-run, the axiom
re-print, this note's claim check \texttt{claim-check.md} --- added to
\texttt{audit/} after the audit department returns its adjudication, and a
prerequisite for publication.
\end{itemize}

\section{Novelty evidence}\label{sec:novelty}

Adjudication copied from the pipeline's exclusivity reviews (C9 gate), which
ran before the problem was accepted: verdict \textbf{no occupying record
found} for the period-3 pendant-ladder matching sequence, its order-12
sparse recurrence, and its three mod-3 subsequences --- at both the object
level (miner screen) and the proposition level (independent full review),
confirmed by a second review that deliberately reused no query method of the
first. Value tier \textbf{new sequence / new recurrence}. The registered
family members are delimited explicitly and none of them is this object:
the plain ladder A030186 \cite{a030186}, the period-1 pendant corona
A102436 \cite{a102436}, the period-2 pendant A386889 \cite{a386889}; the
nearest literature neighbour, Farrell's ``Matchings in ladders''
\cite{farrell1978}, covers the plain ladder only. The claim is strictly
\emph{new sequence / new recurrence}, never new mathematics. Summary by
channel; the full query log, including every zero-hit query, is in this
deposit at \texttt{audit/c9-record.md} (two records, byte-for-byte).

\begin{center}
\resizebox{\textwidth}{!}{%
\begin{tabular}{lll}
\hline
channel & queries & result \\
\hline
OEIS, numeric \cite{oeis} & first review 16 (4 prefix windows, x2/+1/-1,
three subsequences x 3 forms) + second review 8 (shifted-prefix a(0)=1/0,
mid/deep/tail windows, subsequence late windows) & all zero hits;
positive-control anchors hit (Fibonacci, A180970, family keyword sweep) \\
OEIS, recurrence/form & first review 9 (recurrence, coefficients, generating
denominators, characteristic polynomial in two typings) + second review 4
(OEIS-native asterisk typings) & all zero hits \\
OEIS, keyword/alias & first review 12 + second review 8 (faulty/defective
ladder, dimer-monomer phrasing, independent edge sets) + 2 site-scoped web
strings (Chinese layer, site:oeis) &
noise hits only, none this object; reverse cross-reference hunting: A386889
has no reverse references, the A030186 citing family has no period-3 member \\
MSE in-site + MathOverflow & first review: MSE API rate-limited, 2
site-scoped web-search rounds; second review: 5 in-site strings +
MathOverflow (1 query + 2 relaxations) & controls hit (seating-couples
ladder thread); this object absent \\
zbMATH & first review 5 + second review 5 + full-record read & zero hits
for the object; Farrell 1978
(an:0414.05042) read in full: plain ladder, no pendant construction \\
arXiv API & first review 6 + second review 4, with controls & controls hit;
nearest neighbours are graceful
labelling, spin-ladder physics, polymer chemistry --- none this object \\
OpenAlex & 8 incl.\ title-filtered & no this-object title in any sweep \\
library layer & mathlib rg, compfiles, sequencelib tree (26336 files) &
zero graph-ladder or pendant content; mathlib has matching predicates but
no counting API \\
independent re-count & fourth implementation (column-transfer enumeration) &
39/39 agreement with the deep-fit log; guards against single-implementation
blind spots \\
\hline
\end{tabular}%
}
\end{center}

The channels are named so a later reader can re-run them; if any channel
later returns a same-form record, the tier above weakens accordingly, and
this deposit will be amended by a later version rather than by editing this
record.

\section{Scope and limitations}\label{sec:limits}

Grade and boundary, stated as the audit states them. (a) The kernel layer
proves the pure-sequence theorems
Theorems~\ref{thm:interleave}--\ref{thm:mod4} and nothing else. (b) The
combinatorial identification of \texttt{ap3} with the matching count of the
period-3 pendant ladder is \emph{not} kernel-proved: it rests on the
pipeline's phase-0 computational evidence (five methodologically distinct
counting implementations agreeing over 241 terms; double-window modular
scans; rational-domain least-order fitting), is carried at the conjecture /
report layer, and does not become a theorem by this deposit. (c) The
minimality of the periods 15 and 30 is numerically observed (first-return
state scan, unique minimum within 20000 terms) but not kernel-proved here;
the theorems assert the residue laws for the tables as written. (d) Two
further modular laws (period 39 mod 3, period 60 mod 8) and a pure-period
existence statement for every modulus are numerically supported but not
kernel-proved and are not claimed. (e) The perfect-matching (PM)
subsequence has zero-interleaved, too-short non-zero segments for effective
sequence search; it was not novelty-searched and is not claimed. (f) The
gate-four sign-off of the underlying task was exercised by the author's
instruction of 2026-09-29 to run both the paper lane and this claim lane
with direct outbound authorisation, while the task report file header
retains its pre-sign-off wording unchanged; both calibers are recorded.
(g) This note carries no literature review and is a deposit record, not a
survey; the narrative paper is prepared separately under the pipeline's
paper lane and will reference this deposit.

\section*{Data availability}

The deposit package for this claim contains the claim note (LaTeX source and
PDF), the Lean sources (\texttt{proofs/}: frozen statement file, final proof
file, \texttt{lean-toolchain}), the verification and audit logs, and the
full C9 query log with zero-hit entries (\texttt{audit/}). Public mirror:
the claim lane's repository \url{https://github.com/Aurora0134/ai4math-claims},
directory \texttt{ladder-pend3-interleave/} (pushed under the author's
instruction of 2026-09-29; the commit and timestamp are recorded in the
source repository's \texttt{claims/ladder-pend3-interleave/card.md}). Zenodo:
the DOI is reserved at upload by the author and is recorded in
\texttt{metadata/README.md} of this package once created; see
\texttt{UPLOAD.md} in the source repository. Claim note under CC BY 4.0,
code under MIT.

\section*{Use of generative AI}

(a) AI performed: candidate generation and the two-stage exclusivity review,
autoformalisation of the definitions and statements, the roundtrip semantic
check, proof search and proof-script writing (worker agents, with
main-agent repair on one theorem after a tactic-level rejection), error
classification and repair, and the assembly of this note including its
verbatim listings. (b) Model, node and budget: the proving phase ran on the
node designated by the author for that session (\texttt{a6api/gpt-6-sol};
the endpoint/model detection scripts resolved unknown for that session and
the designation is recorded in the task ledger's override log), with no
cross-node switching, consuming $7$ agent-run samples of a permitted $20$,
$0$ llm-calls of a permitted $8$, and $27$ strict kernel compiles of a
permitted $48$ (plus compile-debug probes, separately ledgered), copied
verbatim from the task ledger. This note consumed $0$ pipeline LLM calls
against a cap of $4$; drafting was mechanical assembly. (c) Final
adjudication: all statements and proofs reported here were checked by the
Lean 4 kernel: the proof compiles with no \texttt{sorry}, and
\texttt{\#print axioms} reports only \{propext, Quot.sound,
Classical.choice\}. (d) The AI system is not an author; the human author
reviewed the evidence and is responsible for all content.

\section*{Author contributions}

CRediT-style: the human author adjudicated topic selection, statement
semantics and the deposit decision; the AI4Math pipeline (an LLM) generated
the candidate propositions, the proof searches and the text drafts. The AI
system is not listed as an author.

\begin{thebibliography}{9}
% Each entry verified on 2026-09-29 against Crossref title/page metadata
% (lean4, mathlib), the zbMATH Open API record an:0414.05042 (Farrell), and
% live fetches of oeis.org (A030186, A102436, A386889). Raw responses:
% claims/ladder-pend3-interleave/audit/bib-verify/. No entry is quoted from
% memory.
\bibitem{lean4}
L.~de Moura and S.~Ullrich.
\newblock The Lean 4 theorem prover and programming language.
\newblock \emph{Automated Deduction -- CADE-28}, LNCS, pp.~625--635, 2021.
\url{https://doi.org/10.1007/978-3-030-79876-5_37}

\bibitem{mathlib}
The mathlib Community.
\newblock The Lean mathematical library.
\newblock \emph{CPP 2020}, pp.~367--381, 2020.
\url{https://doi.org/10.1145/3372885.3373824}

\bibitem{oeis}
The OEIS Foundation.
\newblock \emph{The On-Line Encyclopedia of Integer Sequences}.
\newblock \url{https://oeis.org/}

\bibitem{a030186}
OEIS Foundation.
\newblock Matchings of the ladder graph $P_2\times P_n$ (\%N line gives the
defining recurrence $a(n)=3a(n-1)+a(n-2)-a(n-3)$, $a(0)=1$, $a(1)=2$,
$a(2)=7$).
\newblock \emph{The On-Line Encyclopedia of Integer Sequences}, entry
A030186.
\url{https://oeis.org/A030186}

\bibitem{a102436}
OEIS Foundation.
\newblock Number of matchings of the corona $L'(n)$ of the ladder graph
$L(n)=P_2\times P_n$ and the complete graph $K(1)$ (the period-1 pendant
relative).
\newblock \emph{The On-Line Encyclopedia of Integer Sequences}, entry
A102436.
\url{https://oeis.org/A102436}

\bibitem{a386889}
OEIS Foundation.
\newblock Number of ways to use $1\times 1$ squares and $1\times 2$ dominos
(in any orientation) to tile a $3\times(2n-1)$ strip with every other cell
in the top row removed (the period-2 pendant relative).
\newblock \emph{The On-Line Encyclopedia of Integer Sequences}, entry
A386889.
\url{https://oeis.org/A386889}

\bibitem{farrell1978}
E.~J. Farrell.
\newblock Matchings in ladders.
\newblock \emph{Ars Combinatoria} 6 (1978), pp.~153--161.
\newblock (zbMATH \texttt{an:0414.05042}.)

\end{thebibliography}

\end{document}
